% Materiał skalarny: AM I-11, zachowane etykiety wyników.
\chapter{Iloczyny Cauchy'ego}
\label{ma1:ch:cauchy-product}
\noindent Wykorzystujemy wyniki rozdziałów
\ref{ma2:ch:banach-series} i \ref{ma2:ch:unconditional-series}.

Mnożenie dwóch skończonych sum prowadzi do sumy iloczynów wszystkich par składników. Dla szeregów nieskończonych pojawia się problem: w jakiej kolejności należy dodawać te iloczyny i kiedy otrzymana suma jest iloczynem sum wyjściowych? Najpierw rozstrzygniemy to dla szeregów rzeczywistych i zespolonych, a następnie dla ciągłych odwzorowań dwuliniowych między przestrzeniami unormowanymi. Nie korzystamy z teorii funkcji wykładniczej ani szeregów potęgowych.

\section{Definicja i geometria indeksów}
\begin{definicja}
\label{ma1:def:cauchy-product}
Niech $K=\R$ lub $K=\C$, a $(a_n)_{n\ge0}$ i $(b_n)_{n\ge0}$ będą ciągami w $K$. Iloczynem Cauchy'ego szeregów $\sum_{n=0}^\infty a_n$ i $\sum_{n=0}^\infty b_n$ nazywamy szereg $\sum_{n=0}^\infty c_n$, gdzie dla każdego całkowitego $n\ge0$
\[
 c_n=\sum_{j=0}^n a_jb_{n-j}.
\]
Definicja nie wymaga zbieżności żadnego z trzech szeregów.
\end{definicja}
W całym rozdziale przez $\N_0$ oznaczamy zbiór $\{0,1,2,\ldots\}$. Indeksowanie od zera upraszcza wzory: wyraz $c_n$ zawiera dokładnie te iloczyny $a_jb_k$, dla których $j+k=n$. Dla szeregów indeksowanych od jedynki można najpierw przeindeksować wyrazy, przyjmując $\widetilde a_n=a_{n+1}$ i $\widetilde b_n=b_{n+1}$ dla $n\in\N_0$. Takie przeindeksowanie nie zmienia sum, ponieważ $\sum_{n=0}^N\widetilde a_n=\sum_{j=1}^{N+1}a_j$ i analogicznie dla drugiego szeregu.

\begin{lemat}[Tożsamości skończone]
\label{ma1:lem:cauchy-finite-identities}
Niech $(a_n),(b_n)$ będą ciągami w $K=\R$ lub $K=\C$, indeksowanymi od zera, i niech $(c_n)$ będzie określony jak w definicji \ref{ma1:def:cauchy-product}. Dla każdego całkowitego $N\ge0$ oznaczmy
$A_N=\sum_{j=0}^N a_j$, $B_N=\sum_{k=0}^N b_k$ i $C_N=\sum_{n=0}^N c_n$. Wówczas
\[
 C_N=\sum_{\substack{j,k\ge0\\j+k\le N}}a_jb_k
 =\sum_{j=0}^N a_jB_{N-j},
 \qquad
 A_NB_N=\sum_{j=0}^N\sum_{k=0}^N a_jb_k.
\]
We wszystkich sumach indeksy są całkowite.
\end{lemat}
\begin{proof}
Rozwijamy $C_N=\sum_{n=0}^N\sum_{j=0}^n a_jb_{n-j}$. Przypisanie $(n,j)\mapsto(j,n-j)$ jest bijekcją zbioru par całkowitych $0\le j\le n\le N$ na zbiór par nieujemnych $j,k$ z $j+k\le N$. Przypisaniem odwrotnym jest $(j,k)\mapsto(j+k,j)$. Dla ustalonego $j$ dopuszczalne wartości $k$ to $0,\ldots,N-j$, co daje drugą równość. Ostatnia równość wynika z rozdzielności mnożenia względem dodawania. Wszystkie przekształcenia dotyczą sum skończonych.
\end{proof}

\begin{figure}[htbp]
\centering
\begin{tikzpicture}[x=.8cm,y=.8cm,>=Stealth]
\fill[blue!10] (-.22,-.22) rectangle (4.22,4.22);
\fill[teal!22] (-.22,-.22)--(4.22,-.22)--(4.22,.22)--(.22,4.22)--(-.22,4.22)--cycle;
\draw[->] (-.4,0)--(5.3,0) node[right] {$j$};
\draw[->] (0,-.4)--(0,5.3) node[above] {$k$};
\foreach \j in {0,...,4}{
 \foreach \k in {0,...,4}{
  \pgfmathtruncatemacro{\d}{\j+\k}
  \ifnum\d>4
   \fill[blue!65!black] (\j,\k) circle (2pt);
  \else
   \fill[teal!75!black] (\j,\k) circle (2pt);
  \fi
 }
}
\draw[teal!75!black,thick] (0,4)--(4,0);
\draw[densely dashed,blue!65!black] (0,4)--(4,4)--(4,0);
\node[left] at (0,4) {$N$};
\node[below] at (4,0) {$N$};
\node[below left] at (0,0) {$0$};
\node[anchor=west,align=left] at (5.5,3.6) {\small Trójkąt: $j+k\le N$\\\small suma $C_N$};
\node[anchor=west,align=left] at (5.5,1.6) {\small Poza trójkątem: $j+k>N$\\\small część kwadratu $j,k\le N$};
\fill[teal!75!black] (5.25,3.65) circle (2pt);
\fill[blue!65!black] (5.25,1.65) circle (2pt);
\end{tikzpicture}
\caption{Dla $N=4$ suma częściowa iloczynu Cauchy'ego obejmuje zielone punkty trójkąta, a iloczyn sum częściowych wszystkie zielone i niebieskie punkty kwadratu. Odcinki i tło służą jedynie wyróżnieniu skończonych zbiorów indeksów.}
\label{ma1:fig:cauchy-index-sets}
\end{figure}
\FloatBarrier

\section{Iloczyn szeregów bezwzględnie zbieżnych}
\begin{twierdzenie}
\label{ma1:thm:absolute-cauchy-product}
Niech $(a_n),(b_n)$ będą ciągami w $\R$ lub $\C$, indeksowanymi od zera, i niech oba szeregi $\sum_{n=0}^\infty a_n$, $\sum_{n=0}^\infty b_n$ będą bezwzględnie zbieżne do odpowiednio $A,B$. Jeżeli $c_n=\sum_{j=0}^n a_jb_{n-j}$ dla $n\ge0$, to $\sum_{n=0}^\infty c_n$ jest bezwzględnie zbieżny, ma sumę $AB$ oraz
\[
 \sum_{n=0}^\infty |c_n|
 \le \left(\sum_{n=0}^\infty |a_n|\right)
      \left(\sum_{n=0}^\infty |b_n|\right).
\]
\end{twierdzenie}
\begin{proof}
Oznaczmy $\alpha=\sum_{n=0}^\infty|a_n|$, $\beta=\sum_{n=0}^\infty|b_n|$. Dla każdego całkowitego $N\ge0$ z nierówności trójkąta, lematu \ref{ma1:lem:cauchy-finite-identities} zastosowanego do modułów oraz nieujemności składników mamy
\[
 \sum_{n=0}^N|c_n|\le
 \sum_{\substack{j,k\ge0\\j+k\le N}}|a_j||b_k|
 \le \sum_{j=0}^N\sum_{k=0}^N |a_j||b_k|
 \le\alpha\beta.
\]
Sumy częściowe szeregu modułów są niemalejące i ograniczone z góry przez $\alpha\beta$. Z kryterium zbieżności szeregów nieujemnych (tom I, rozdział 09) szereg modułów jest zbieżny, a jego suma spełnia podane oszacowanie.

Niech $A_N,B_N,C_N$ oznaczają sumy częściowe z lematu \ref{ma1:lem:cauchy-finite-identities}. Z tożsamości skończonych wynika $A_NB_N-C_N=\sum_{0\le j,k\le N,\ j+k>N}a_jb_k$. Dla $m=\lfloor N/2\rfloor$, jeżeli $0\le j,k\le N$ i $j+k>N$, to $j>m$ lub $k>m$: w przeciwnym razie $j+k\le2m\le N$. Zatem
\[
 |A_NB_N-C_N|
 \le \beta\sum_{j>m}|a_j|+
      \alpha\sum_{k>m}|b_k|.
\]
Istotnie, sumę modułów po zbiorze brakujących par ograniczamy przez sumę po dwóch zbiorach: $m<j\le N$, $0\le k\le N$ oraz $0\le j\le N$, $m<k\le N$. Wspólne pary można policzyć dwukrotnie, gdyż składniki są nieujemne. Sumy po tych zbiorach rozkładają się na iloczyny i są nie większe od podanych oszacowań. Gdy $N\to\infty$, mamy $m\to\infty$, a obie reszty szeregów modułów dążą do zera, ponieważ $\sum_{j>m}|a_j|=\alpha-\sum_{j=0}^{m}|a_j|\to0$, a dla drugiego szeregu zachodzi analogiczna równość. Oszacowanie daje $|A_NB_N-C_N|\to0$. Ponieważ $A_NB_N\to AB$, z nierówności $|C_N-AB|\le |C_N-A_NB_N|+|A_NB_N-AB|$ otrzymujemy $C_N\to AB$.
\end{proof}

\begin{przyklad}
\label{ma1:ex:geometric-product}
Niech $q\in\C$ i $|q|<1$. Biorąc $a_n=b_n=q^n$ dla $n\ge0$, otrzymujemy $c_n=\sum_{j=0}^n q^jq^{n-j}=(n+1)q^n$ (dla $q=0$ stosujemy konwencję $q^0=1$). Oba szeregi geometryczne są bezwzględnie zbieżne, ponieważ $|q^n|=|q|^n$ i $|q|<1$. Z twierdzenia \ref{ma2:thm:vector-geometric}, zastosowanego do $v=1$, każdy ma sumę $1/(1-q)$, więc
\[
 \sum_{n=0}^\infty (n+1)q^n=\frac1{(1-q)^2}.
\]
Dla $q=1/2$ suma wynosi $4$. Wzór wynika z mnożenia szeregów, bez różniczkowania i bez funkcji potęgowych o wykładnikach rzeczywistych.
\end{przyklad}

\section{Twierdzenie Mertensa}
\begin{twierdzenie}[Mertensa]
\label{ma1:thm:mertens-product}
Niech $(a_n),(b_n)$ będą ciągami w $\R$ lub $\C$, indeksowanymi od zera. Załóżmy, że $\sum_{n=0}^\infty a_n=A$, $\sum_{n=0}^\infty b_n=B$ i co najmniej jeden z tych szeregów jest zbieżny bezwzględnie. Wtedy iloczyn Cauchy'ego, czyli szereg o wyrazach $c_n=\sum_{j=0}^n a_jb_{n-j}$, jest zbieżny do $AB$.
\end{twierdzenie}
\begin{proof}
Możemy założyć, że szereg $\sum a_n$ jest bezwzględnie zbieżny. Istotnie, dla każdego $n\in\N_0$ podstawienie $k=n-j$ w skończonej sumie oraz przemienność mnożenia dają $\sum_{j=0}^n a_jb_{n-j}=\sum_{k=0}^n b_ka_{n-k}$, więc zamiana ról obu ciągów nie zmienia $c_n$. Oznaczmy $\alpha=\sum_{n=0}^\infty|a_n|$ oraz dla każdego całkowitego $N\ge0$ przyjmijmy $A_N=\sum_{j=0}^N a_j$, $B_N=\sum_{k=0}^N b_k$, $C_N=\sum_{n=0}^N c_n$. Z $B_n\to B$ wynika ograniczoność ciągu $(B_n-B)$, więc istnieje $M\ge1$ takie, że $|B_n-B|\le M$ dla każdego $n\ge0$.

Z lematu \ref{ma1:lem:cauchy-finite-identities} otrzymujemy
\[
 C_N-A_N B=\sum_{j=0}^N a_j(B_{N-j}-B).
\]
Ustalmy $\varepsilon>0$. Wybierzmy $J\in\N_0$ z $M\sum_{j>J}|a_j|<\varepsilon/2$, korzystając ze zbieżności szeregu modułów i ustalonego skończonego $M$. Ponieważ $B_n\to B$, istnieje $L\ge0$ takie, że $|B_n-B|<\varepsilon/(2(1+\alpha))$ dla $n\ge L$. Dla $N\ge J+L$ oraz $0\le j\le J$ mamy $N-j\ge L$, zatem
\[
\begin{aligned}
 |C_N-A_N B|
 &\le\sum_{j=0}^J |a_j||B_{N-j}-B|
      +\sum_{j=J+1}^N |a_j||B_{N-j}-B|\\
 &\le\frac{\varepsilon\alpha}{2(1+\alpha)}
        +M\sum_{j>J}|a_j|
 <\varepsilon.
\end{aligned}
\]
Druga suma jest zerowa, jeśli $N=J$. Wobec tego $C_N-A_NB\to0$. Ponieważ $A_N\to A$, otrzymujemy $C_N\to AB$.
\end{proof}

\begin{przyklad}
Niech $a_n=2^{-n}$ oraz $b_n=(-1)^n/(n+1)$ dla $n\ge0$. Pierwszy szereg ma sumę $2$ i jest zbieżny bezwzględnie, drugi jest zbieżny warunkowo z kryterium Leibniza. Jeżeli $S=\sum_{n=0}^\infty(-1)^n/(n+1)$, twierdzenie Mertensa daje
\[
 \sum_{n=0}^\infty\left(\sum_{j=0}^n\frac{(-1)^{n-j}}{2^j(n-j+1)}\right)=2S.
\]
Iloczyn ten jest zbieżny warunkowo. Istotnie, z definicji $c_0=b_0$ i dla $n\ge1$ zachodzi $c_n=b_n+\frac12c_{n-1}$ (wydzielamy składnik $j=0$ i w pozostałej sumie zmieniamy indeks na $j-1$). Gdyby $\sum_{n=0}^\infty|c_n|$ był zbieżny, to szereg $\sum_{n=1}^\infty(|c_n|+\frac12|c_{n-1}|)$ również byłby zbieżny na mocy liniowości i niezmienniczości zbieżności po pominięciu skończonego początku. Z $|b_n|\le |c_n|+\frac12|c_{n-1}|$ dla $n\ge1$ i kryterium porównawczego wynikałaby zbieżność $\sum_{n=0}^\infty|b_n|$, sprzecznie z rozbieżnością szeregu harmonicznego.
\end{przyklad}

\section{Iloczyn dwóch szeregów zbieżnych warunkowo}
\begin{przyklad}
\label{ma1:ex:divergent-cauchy-product}
Niech $a_n=b_n=(-1)^n/\sqrt{n+1}$ dla całkowitych $n\ge0$. Oba szeregi są zbieżne z kryterium Leibniza: $1/\sqrt{n+1}$ jest ciągiem malejącym, a dla $\varepsilon>0$ nierówność $n+1>1/\varepsilon^2$ daje $1/\sqrt{n+1}<\varepsilon$. Nie są zbieżne bezwzględnie, ponieważ $1/\sqrt{n+1}\ge1/(n+1)$, a szereg harmoniczny jest rozbieżny.
Jednak
\[
 c_n=(-1)^n\sum_{j=0}^n
       \frac1{\sqrt{(j+1)(n-j+1)}}.
\]
Dla $0\le j\le n$ mamy $j+1\le n+1$ i $n-j+1\le n+1$, zatem $\sqrt{(j+1)(n-j+1)}\le n+1$. Wszystkie składniki sumy bez znaku są dodatnie, a jest ich $n+1$. Multiplikatywność nieujemnych pierwiastków kwadratowych wynika z jednoznaczności pierwiastka: iloczyn dwóch nieujemnych pierwiastków jest nieujemny, a jego kwadrat jest iloczynem liczb podpierwiastkowych. Wobec tego
\[
 |c_n|\ge(n+1)\frac1{n+1}=1.
\]
Wyrazy iloczynu Cauchy'ego nie dążą do zera. Iloczyn ten jest rozbieżny na mocy twierdzenia \ref{ma2:thm:series-cauchy}.
\end{przyklad}
Pytanie o istnienie sumy iloczynu Cauchy'ego trzeba odróżnić od pytania o jej wartość. Następne twierdzenie pokazuje, że jeśli wszystkie trzy szeregi są zbieżne, to suma iloczynu musi być iloczynem sum, nawet bez zbieżności bezwzględnej.

\section{Twierdzenie Abela o iloczynie Cauchy'ego}
\begin{lemat}
\label{ma1:lem:triangular-means}
Niech $K=\R$ lub $K=\C$, a $(x_n)_{n\ge0},(y_n)_{n\ge0}$ będą ciągami w $K$ zbieżnymi odpowiednio do $x,y\in K$. Wówczas
\[
 \lim_{N\to\infty}\frac1{N+1}\sum_{j=0}^N x_jy_{N-j}=xy.
\]
\end{lemat}
\begin{proof}
Ciąg $(y_n)$ jest ograniczony, więc wybieramy $M\ge1$ z $|y_n|\le M$ dla każdego $n$ oraz $|x|\le M$. Dla dowolnych całkowitych $j,k\ge0$ rozkład $x_jy_k-xy=(x_j-x)y_k+x(y_k-y)$ i nierówność trójkąta dają
$|x_jy_k-xy|\le M(|x_j-x|+|y_k-y|)$.
Ustalmy $\varepsilon>0$. Istnieje całkowite $J\ge1$ takie, że $|x_j-x|<\varepsilon/(4M)$ i $|y_j-y|<\varepsilon/(4M)$ dla każdego $j\ge J$. Jeżeli $N\ge2J$, to dla $J\le j\le N-J$ zachodzi $|x_jy_{N-j}-xy|<\varepsilon/2$.

Z ograniczoności obu ciągów istnieje rzeczywiste $D>0$ takie, że $|x_jy_k-xy|\le D$ dla wszystkich $j,k\ge0$. Na przykład dla ograniczeń $|x_j|\le P$, $|y_k|\le Q$ można wziąć $D=PQ+|xy|+1$. Dla $N\ge2J$ w sumie po $0\le j\le N$ dokładnie $2J$ indeksów nie spełnia $J\le j\le N-J$: po $J$ indeksów przy każdym końcu. W konsekwencji
\[
 \left|\frac1{N+1}\sum_{j=0}^N x_jy_{N-j}-xy\right|
 \le \frac{\varepsilon}{2}+\frac{2JD}{N+1}.
\]
Dla każdego całkowitego $N\ge2J$ spełniającego ponadto $N+1>4JD/\varepsilon$ drugi składnik jest mniejszy od $\varepsilon/2$. Taki próg istnieje z własności Archimedesa, co dowodzi tezy.
\end{proof}
\begin{twierdzenie}[Abela o iloczynie Cauchy'ego]
\label{ma1:thm:abel-cauchy-product}
Niech $(a_n),(b_n)$ będą ciągami liczb rzeczywistych lub zespolonych, indeksowanymi od zera. Jeżeli szeregi $\sum_{n=0}^\infty a_n$, $\sum_{n=0}^\infty b_n$ oraz ich iloczyn Cauchy'ego $\sum_{n=0}^\infty c_n$, gdzie $c_n=\sum_{j=0}^n a_jb_{n-j}$, są zbieżne odpowiednio do $A,B,C$, to $C=AB$.
\end{twierdzenie}
\begin{proof}
Oznaczmy sumy częściowe przez $A_N,B_N,C_N$ jak w lemacie \ref{ma1:lem:cauchy-finite-identities}. Dla każdego całkowitego $N\ge0$ zachodzi tożsamość
\[
 \sum_{m=0}^N C_m=\sum_{j=0}^N A_j B_{N-j}.
\]
Aby ją uzasadnić, ustalmy parę nieujemnych indeksów całkowitych $r,s$ z $r+s\le N$. Składnik $a_rb_s$ po lewej występuje w sumach $C_m$ dokładnie dla $r+s\le m\le N$, a więc $N-r-s+1$ razy. Po prawej występuje dokładnie dla $r\le j\le N-s$, również $N-r-s+1$ razy. Jeżeli $r+s>N$, nie występuje po żadnej ze stron. Ponieważ rozpatrujemy skończone sumy, równość liczności wystąpień dla każdej pary dowodzi tożsamości.

Dzielimy obie strony przez $N+1$. Z lematu \ref{ma1:lem:triangular-means} zastosowanego do $x_j=A_j$, $y_j=B_j$ wynika, że prawa strona po podzieleniu dąży do $AB$. Stosując ten sam lemat do $x_j=C_j$ oraz stałego ciągu $y_j=1$, otrzymujemy, że lewa strona po podzieleniu dąży do $C$. Z jednoznaczności granicy mamy $C=AB$.
\end{proof}
Zaletą tego dowodu jest niezależność od teorii szeregów potęgowych.

\begin{przyklad}
\label{ma1:ex:convergent-conditional-product}
Niech $a_n=b_n=(-1)^n/(n+1)$ dla $n\in\N_0$. Oba szeregi są warunkowo zbieżne do tej samej sumy $S$ (przykład \ref{ma1:ex:alternating-harmonic}). Tym razem ich iloczyn Cauchy'ego jest zbieżny. Dla każdego $n\ge0$ mamy
\[
 c_n=(-1)^n\sum_{j=0}^n\frac1{(j+1)(n-j+1)}
     =(-1)^n\frac{2H_{n+1}}{n+2},
 \qquad H_m=\sum_{j=1}^m\frac1j\quad(m\ge1).
\]
Druga równość wynika z tożsamości
$\frac1{(j+1)(n-j+1)}=\frac1{n+2}(\frac1{j+1}+\frac1{n-j+1})$ i zmiany indeksu w drugiej sumie.
Ciąg $H_m/(m+1)$ jest nierosnący, ponieważ
\[
 \frac{H_m}{m+1}-\frac{H_{m+1}}{m+2}
 =\frac{H_m-1}{(m+1)(m+2)}\ge0.
\]
Ponadto dąży do zera: dla $\varepsilon>0$ wybieramy $J\ge1$ z $1/(J+1)<\varepsilon/2$. Dla $m\ge J$ mamy $H_m\le H_J+(m-J)/(J+1)$, stąd $0\le H_m/(m+1)\le H_J/(m+1)+1/(J+1)<\varepsilon$ dla dostatecznie dużego $m$.
Z kryterium Leibniza szereg $\sum_{n=0}^\infty c_n$ jest zbieżny, a z twierdzenia \ref{ma1:thm:abel-cauchy-product} jego suma wynosi $S^2$. Nie jest bezwzględnie zbieżny, gdyż $2H_{n+1}/(n+2)\ge2/(n+2)$, a szereg po prawej jest rozbieżny przez porównanie z szeregiem harmonicznym po pominięciu pierwszego wyrazu. Wraz z przykładem \ref{ma1:ex:divergent-cauchy-product} pokazuje to, że sama warunkowa zbieżność obu czynników nie rozstrzyga zbieżności iloczynu.
\end{przyklad}
 Wykorzystana własność średnich arytmetycznych pojawi się ponownie przy omawianiu metod sumowania szeregów.

\section{Sumowanie iloczynów po wszystkich parach indeksów}
\begin{twierdzenie}
\label{ma1:thm:pair-enumeration}
Niech $(a_n)_{n\ge0},(b_n)_{n\ge0}$ będą ciągami w $\R$ lub $\C$ takimi, że oba szeregi są zbieżne bezwzględnie, odpowiednio do $A,B$. Dla dowolnej bijekcji $\pi:\N_+\to\N_0\times\N_0$, gdzie $\pi(l)=(j_l,k_l)$, szereg
\[
 \sum_{l=1}^\infty a_{j_l}b_{k_l}
\]
jest zbieżny bezwzględnie i ma sumę $AB$.
\end{twierdzenie}
\begin{proof}
Istnieje wyliczenie wszystkich par w kolejności rosnącej sumy współrzędnych: dla kolejnych $n=0,1,\ldots$ wypisujemy $(0,n),(1,n-1),\ldots,(n,0)$. Każda para $(j,k)$ należy do dokładnie jednego bloku, o numerze $n=j+k$, i występuje w nim dokładnie raz. Blok o numerze $n$ ma $n+1$ elementów, a przed nim jest $\sum_{i=0}^{n-1}(i+1)=n(n+1)/2$ elementów (suma pusta dla $n=0$). Wobec tego pozycja pary $(j,k)$, dla $n=j+k$, jest równa $n(n+1)/2+j+1$. Bloki wypełniają kolejne rozłączne przedziały dodatnich indeksów, a ich długości mają nieograniczone sumy częściowe. Otrzymujemy bijekcję $\pi_0:\N_+\to\N_0\times\N_0$.
Oznaczmy $\alpha=\sum_{n\ge0}|a_n|$ i $\beta=\sum_{n\ge0}|b_n|$. Dla dowolnego skończonego początkowego fragmentu tego wyliczenia istnieje całkowite $M\ge0$ takie, że wszystkie jego pary mają obie współrzędne nie większe od $M$. Suma modułów odpowiadających tym parom jest nie większa od
$(\sum_{j=0}^M|a_j|)(\sum_{k=0}^M|b_k|)\le\alpha\beta$. Zatem szereg wyrazów $a_jb_k$ wyliczonych przez $\pi_0$ jest zbieżny bezwzględnie, a więc zbieżny.
Grupując jego kolejne bloki według $j+k=n$, otrzymujemy iloczyn Cauchy'ego. Twierdzenie \ref{ma2:thm:series-grouping} zapewnia zachowanie sumy, a twierdzenie \ref{ma1:thm:absolute-cauchy-product} identyfikuje ją z $AB$. Dowolna bijekcja $\pi$ wyznacza przestawienie wyliczenia $\pi_0$, mianowicie bijekcję $\pi_0^{-1}\circ\pi$ zbioru $\N_+$ na siebie. Z twierdzenia \ref{ma2:thm:absolute-rearrangement} każde takie przestawienie jest zbieżne bezwzględnie do tej samej sumy.
\end{proof}

\section{Wielokrotny iloczyn Cauchy'ego}
\begin{wniosek}
\label{ma1:cor:finite-cauchy-products}
Niech $r\in\N_+$ i niech dla każdego $i\in\{1,\ldots,r\}$ szereg $\sum_{n=0}^\infty a_n^{(i)}$ w $\R$ lub $\C$ będzie zbieżny bezwzględnie do $A_i$. Dla każdego $n\in\N_0$ połóżmy
\[
 d_n=\sum_{\substack{n_1,\ldots,n_r\in\N_0\\
                         n_1+\cdots+n_r=n}}
                    \prod_{i=1}^r a_{n_i}^{(i)}.
\]
Każda suma definiująca $d_n$ jest skończona, ponieważ wszystkie $n_i$ należą do $\{0,\ldots,n\}$. Wówczas szereg $\sum_{n=0}^\infty d_n$ jest zbieżny bezwzględnie i ma sumę $\prod_{i=1}^r A_i$. Ponadto
\[
 \sum_{n=0}^\infty |d_n|
 \le \prod_{i=1}^r\left(\sum_{n=0}^\infty |a_n^{(i)}|\right).
\]
\end{wniosek}
\begin{proof}
Dla $r=1$ teza jest bezpośrednia. Załóżmy ją dla $r$ szeregów i rozważmy $r+1$ szeregów. Oznaczmy współczynniki iloczynu pierwszych $r$ przez $d_j$. Współczynnik iloczynu Cauchy'ego tego szeregu i ostatniego szeregu jest równy
\[
 \sum_{j=0}^n d_j a_{n-j}^{(r+1)}
 =\sum_{\substack{n_1,\ldots,n_{r+1}\in\N_0\\
                        n_1+\cdots+n_{r+1}=n}}
             \prod_{i=1}^{r+1} a_{n_i}^{(i)}.
\]
Istotnie, każda krotka po prawej występuje dokładnie raz po lewej, dla $j=n_1+\cdots+n_r$. Wszystkie te sumy są skończone, bo każdy indeks należy do $\{0,\ldots,n\}$. Z założenia indukcyjnego szereg $\sum d_j$ jest bezwzględnie zbieżny. Stosując twierdzenie \ref{ma1:thm:absolute-cauchy-product}, otrzymujemy bezwzględną zbieżność nowego szeregu oraz iloczyn sum i oszacowanie sumy modułów przez iloczyn poprzedniego oszacowania i $\sum_{n\ge0}|a_n^{(r+1)}|$. Indukcja kończy dowód.
\end{proof}
\begin{uwaga}
Iloczyn Cauchy'ego nie oznacza mnożenia wyraz po wyrazie, czyli szeregu $\sum a_nb_n$. Pozwala on natomiast przejść od iloczynu dwóch szeregów do sumy współczynników wyznaczanych przez skończone sumy. Będzie to podstawą dowodu prawa dodawania argumentów funkcji wykładniczej. Poniżej rozszerzymy te wyniki na odwzorowania dwuliniowe.
\end{uwaga}

\section{Iloczyn Cauchy'ego względem odwzorowania dwuliniowego}
Niech $E,F,G$ będą przestrzeniami unormowanymi nad tym samym
ciałem $\K=\R$ lub $\K=\C$. Nie zakładamy istnienia
mnożenia wektorów w żadnej z tych przestrzeni.

\begin{lemat}[Ciągłość odwzorowania dwuliniowego]
\label{ma2:lem:bilinear-bound}
Dla odwzorowania dwuliniowego $T:E\times F\to G$
następujące warunki są równoważne: ciągłość $T$ względem
normy $\|(x,y)\|=\max(\|x\|,\|y\|)$ oraz istnienie
stałej $K\ge0$ takiej, że dla wszystkich $x\in E$, $y\in F$
\[
 \|T(x,y)\|\le K\|x\|\|y\|.
\]
\end{lemat}
\begin{proof}
Jeżeli $T$ jest ciągłe, to z $T(0,0)=0$ istnieje
$\delta>0$ takie, że $\|T(u,v)\|<1$ dla
$\max(\|u\|,\|v\|)<\delta$. Dla $x\ne0$, $y\ne0$
podstawiamy $u=\delta x/(2\|x\|)$ i $v=\delta y/(2\|y\|)$.
Obie normy są równe $\delta/2$, więc dwuliniowość daje
$\|T(x,y)\|<4\|x\|\|y\|/\delta^2$.
Jeśli którykolwiek wektor jest zerowy, $T(x,y)=0$.
Można zatem przyjąć $K=4/\delta^2$.
Odwrotnie, z dwuliniowości
\[
 T(x',y')-T(x,y)=T(x'-x,y')+T(x,y'-y).
\]
Jeśli $(x',y')\to(x,y)$, to $\|y'\|\le\|y\|+1$ od
pewnego miejsca. Norma różnicy nie przekracza
$K(\|x'-x\|(\|y\|+1)+\|x\|\|y'-y\|)$, co dąży do zera.
Kryterium ciągowe ciągłości w przestrzeniach metrycznych
daje ciągłość w każdym punkcie.
\end{proof}

\begin{definicja}
\label{ma2:def:bilinear-cauchy}
Dla ciągłego odwzorowania dwuliniowego $T:E\times F\to G$
oraz ciągów $(x_n)_{n\ge0}$ w $E$ i $(y_n)_{n\ge0}$ w $F$
określamy iloczyn Cauchy'ego względem $T$ przez
$z_n=\sum_{j=0}^n T(x_j,y_{n-j})$. Jest to ciąg w $G$.
Oznaczamy
\[
 X_N=\sum_{j=0}^N x_j,\qquad Y_N=\sum_{k=0}^N y_k,\qquad
 Z_N=\sum_{n=0}^N z_n.
\]
\end{definicja}
Dwuliniowość oraz bijekcja indeksów z lematu
\ref{ma1:lem:cauchy-finite-identities} dają
\[
\begin{aligned}
 Z_N&=\sum_{j+k\le N}T(x_j,y_k)
 =\sum_{j=0}^N T(x_j,Y_{N-j})
 =\sum_{k=0}^N T(X_{N-k},y_k),\\
 T(X_N,Y_N)&=\sum_{j=0}^N\sum_{k=0}^N T(x_j,y_k).
\end{aligned}
\]
W pierwszej sumie $j,k\in\N_0$; wszystkie sumy są skończone.
Druga postać $Z_N$ wynika z zebrania składników dla
ustalonego $j$ i liniowości w drugim argumencie $T$;
trzecia z analogicznego zebrania składników dla ustalonego
$k$ i liniowości w pierwszym argumencie. Rozwinięcie
$T(X_N,Y_N)$ wykorzystuje kolejno liniowość w obu argumentach.
Nie zamieniamy argumentów $T$, gdyż przestrzenie $E,F$
mogą być różne i przemienność nie jest założona.

\begin{twierdzenie}[Mertensa w wersji dwuliniowej]
\label{ma2:thm:bilinear-mertens}
Niech $T:E\times F\to G$ będzie ciągłe dwuliniowe.
Załóżmy, że $\sum x_n=x\in E$ i $\sum y_n=y\in F$.
Jeżeli co najmniej jeden z szeregów jest bezwzględnie
zbieżny, to $\sum z_n=T(x,y)$ w $G$.
Nie wymaga to zupełności żadnej z przestrzeni.
\end{twierdzenie}
\begin{proof}
Wybieramy $K\ge1$ spełniające oszacowanie z lematu
\ref{ma2:lem:bilinear-bound}. Najpierw załóżmy
$\alpha=\sum_{j\ge0}\|x_j\|<\infty$.
Z $Y_m\to y$ wynika istnienie $M\ge1$ z
$\|Y_m-y\|\le M$ dla każdego $m\ge0$.
Z tożsamości skończonych otrzymujemy
$Z_N-T(X_N,y)=\sum_{j=0}^N T(x_j,Y_{N-j}-y)$.
Dla $\varepsilon>0$ wybieramy $J\ge0$ tak, by
$KM\sum_{j>J}\|x_j\|<\varepsilon/2$.
Następnie wybieramy $L\ge0$ tak, by
$\|Y_m-y\|<\varepsilon/(2K(1+\alpha))$ dla $m\ge L$.
Dla $N\ge J+L$ mamy $N-j\ge L$ przy $j\le J$, więc
\[
\begin{aligned}
 \|Z_N-T(X_N,y)\|
 &\le K\sum_{j=0}^J\|x_j\|\|Y_{N-j}-y\|
    +K\sum_{j=J+1}^N\|x_j\|\|Y_{N-j}-y\|\\
 &\le \frac{\varepsilon\alpha}{2(1+\alpha)}
      +KM\sum_{j>J}\|x_j\|<\varepsilon.
\end{aligned}
\]
Druga suma może być pusta. Ponadto
$\|T(X_N,y)-T(x,y)\|\le K\|X_N-x\|\|y\|\to0$.
Nierówność trójkąta daje $Z_N\to T(x,y)$.

Jeżeli zamiast tego $\sum\|y_k\|<\infty$, stosujemy
już udowodniony przypadek do odwzorowania
$\widetilde T:F\times E\to G$ określonego przez
$\widetilde T(v,u)=T(u,v)$. Jest ono dwuliniowe
i spełnia to samo oszacowanie z $K$, więc jest ciągłe.
Jego współczynnik iloczynu szeregu $(y_n)$ i $(x_n)$
wynosi $\sum_{k=0}^n T(x_{n-k},y_k)=z_n$ po zmianie
indeksu $j=n-k$. Otrzymana suma jest równa
$\widetilde T(y,x)=T(x,y)$, co kończy dowód.
\end{proof}

\begin{wniosek}[Dwa szeregi bezwzględnie zbieżne]
\label{ma2:cor:bilinear-absolute}
Niech $T:E\times F\to G$ będzie ciągłe dwuliniowe,
a $K\ge0$ niech spełnia $\|T(u,v)\|\le K\|u\|\|v\|$
dla wszystkich $u\in E$, $v\in F$. Jeżeli oba szeregi
mają sumy $x,y$ i oba szeregi norm są zbieżne, to iloczyn Cauchy'ego
jest bezwzględnie zbieżny do $T(x,y)$ i
\[
 \sum_{n=0}^{\infty}\|z_n\|
 \le K\left(\sum_{j=0}^\infty\|x_j\|\right)
      \left(\sum_{k=0}^\infty\|y_k\|\right).
\]
W szczególności gdy $E,F$ są przestrzeniami Banacha,
wystarczy założyć zbieżność obu szeregów norm.
Zupełność $G$ nie jest potrzebna, ponieważ granica
iloczynu jest konkretnym elementem $T(x,y)\in G$.
\end{wniosek}
\begin{proof}
Z oszacowania dwuliniowego i nierówności trójkąta
dla każdego $N\ge0$ mamy
\[
 \sum_{n=0}^N\|z_n\|
 \le K\sum_{j+k\le N}\|x_j\|\|y_k\|
 \le K\left(\sum_{j=0}^N\|x_j\|\right)
      \left(\sum_{k=0}^N\|y_k\|\right).
\]
Ostatni iloczyn jest ograniczony przez iloczyn pełnych
sum norm. Sumy częściowe szeregu norm są niemalejące,
więc mają skończoną granicę spełniającą oszacowanie.
Twierdzenie \ref{ma2:thm:bilinear-mertens} daje
zbieżność szeregu wektorowego do $T(x,y)$.
Gdy $E,F$ są zupełne, istnienie $x,y$ wynika
z twierdzenia \ref{ma2:thm:absolute-series} po przeindeksowaniu.
\end{proof}

\begin{lemat}[Średnie po przekątnych]
\label{ma2:lem:bilinear-diagonal-means}
Jeśli $u_j\to u$ w $E$, $v_j\to v$ w $F$ i $T$ jest
ciągłe dwuliniowe, to
\[
 \frac1{N+1}\sum_{j=0}^N T(u_j,v_{N-j})\longrightarrow T(u,v).
\]
\end{lemat}
\begin{proof}
Wybierzmy $K\ge1$ z lematu \ref{ma2:lem:bilinear-bound}
oraz $M\ge1$ z $\|v_k\|,\|u\|\le M$ dla wszystkich $k$.
Dla wszystkich $j,k$ dwuliniowość daje
\[
 \|T(u_j,v_k)-T(u,v)\|
 \le KM(\|u_j-u\|+\|v_k-v\|).
\]
Dla $\varepsilon>0$ wybieramy $J\ge1$ z
$\|u_j-u\|,\|v_j-v\|<\varepsilon/(4KM)$ dla $j\ge J$.
Jeżeli $N\ge2J$ i $J\le j\le N-J$, norma różnicy
jest mniejsza od $\varepsilon/2$.
Z ograniczoności obu ciągów istnieją $P,Q\ge1$
z $\|u_j\|\le P$, $\|v_k\|\le Q$. Dla wszystkich par
norma różnicy nie przekracza $D=KPQ+\|T(u,v)\|$.
Poza środkowym przedziałem jest dokładnie $2J$ indeksów,
więc norma różnicy średniej od $T(u,v)$ jest nie większa
niż $\varepsilon/2+2JD/(N+1)$. Drugi składnik
jest mniejszy od $\varepsilon/2$ dla dostatecznie dużego $N$.
\end{proof}

\begin{twierdzenie}[Abela w wersji dwuliniowej]
\label{ma2:thm:bilinear-abel}
Niech $T:E\times F\to G$ będzie ciągłe dwuliniowe,
a $(z_n)$ niech będzie ciągiem współczynników iloczynu
Cauchy'ego z definicji \ref{ma2:def:bilinear-cauchy}.
Jeżeli $\sum x_n=x$, $\sum y_n=y$ oraz $\sum z_n=z$
są zbieżne w odpowiednich przestrzeniach, to $z=T(x,y)$.
Nie potrzeba zupełności ani zbieżności bezwzględnej.
\end{twierdzenie}
\begin{proof}
Dla $N\ge0$ zachodzi
$\sum_{m=0}^N Z_m=\sum_{j=0}^N T(X_j,Y_{N-j})$.
Istotnie, dla każdej pary $r,s\ge0$ z $r+s\le N$
składnik $T(x_r,y_s)$ pojawia się po lewej w $Z_m$
dla $r+s\le m\le N$, po prawej dla $r\le j\le N-s$,
czyli po obu stronach dokładnie $N-r-s+1$ razy.
Pozostałe pary nie występują. Wszystkie sumy są skończone.
Dzielimy przez $N+1$. Prawa strona dąży do $T(x,y)$
na mocy lematu \ref{ma2:lem:bilinear-diagonal-means}.
Lewa strona dąży do $z$: stosujemy ten sam lemat
do odwzorowania $(w,a)\mapsto aw$ z $G\times\K$ do $G$,
które jest ciągłe dwuliniowe z oszacowaniem $K=1$,
oraz ciągów $Z_j\to z$ i $a_j=1$. Jednoznaczność granicy
daje tezę.
\end{proof}

\begin{przyklad}[Mnożenie macierzy]
\label{ma2:ex:matrix-cauchy-product}
W przestrzeni macierzy kwadratowych $d\times d$, $d\ge1$,
z normą operatorową indukowaną normą euklidesową,
odwzorowanie $T(U,V)=UV$ jest dwuliniowe i ciągłe:
dla każdego wektora $\xi$ mamy
$\|UV\xi\|\le\|U\|\|V\xi\|\le\|U\|\|V\|\|\xi\|$,
więc po wzięciu supremum po $\|\xi\|\le1$ otrzymujemy
$\|UV\|\le\|U\|\|V\|$. Jeśli szeregi macierzy $\sum_{j\ge0}A_j=A$ oraz
$\sum_{j\ge0}B_j=B$ są zbieżne i co najmniej jeden
jest bezwzględnie zbieżny,
twierdzenie \ref{ma2:thm:bilinear-mertens} daje
\[
 \sum_{n=0}^\infty\sum_{j=0}^n A_j B_{n-j}=AB
\]
dla wyrazów szeregów $A_j,B_j$ (nie ich sum częściowych).
W ogólności nie wolno zamienić kolejności mnożenia:
dla $d=2$ i macierzy jednostek $E_{12},E_{21}$ mamy
$E_{12}E_{21}=E_{11}\ne E_{22}=E_{21}E_{12}$, jak wynika
bezpośrednio z definicji mnożenia macierzy.
\end{przyklad}

\begin{wniosek}[Dowolne wyliczenie par w przypadku dwuliniowym]
\label{ma2:cor:bilinear-pair-enumeration}
Niech $T:E\times F\to G$ będzie ciągłe dwuliniowe.
Jeżeli szeregi $\sum x_n=x$ i $\sum y_n=y$ są bezwzględnie
zbieżne oraz $G$ jest przestrzenią Banacha, to dla każdej
bijekcji $\pi:\N_+\to\N_0\times\N_0$,
$\pi(l)=(j_l,k_l)$, szereg
$\sum_{l\ge1}T(x_{j_l},y_{k_l})$ jest bezwzględnie
zbieżny do $T(x,y)$.
\end{wniosek}
\begin{proof}
Wybierzmy $K\ge0$ z lematu \ref{ma2:lem:bilinear-bound}.
Dla dowolnego skończonego początkowego fragmentu wyliczenia
wszystkie pary należą do pewnego kwadratu
$\{0,\ldots,M\}^2$. Z iniektywności wyliczenia i nieujemności
norm otrzymujemy oszacowanie sumy norm tego fragmentu przez
\[
 K\left(\sum_{j=0}^\infty\|x_j\|\right)
  \left(\sum_{k=0}^\infty\|y_k\|\right)<\infty.
\]
Zatem każde takie wyliczenie daje szereg bezwzględnie
zbieżny, a z zupełności $G$ także zbieżny.
Dla wyliczenia po kolejnych przekątnych $j+k=n$,
opisanego w dowodzie twierdzenia
\ref{ma1:thm:pair-enumeration}, grupowanie kolejnych
bloków zachowuje sumę na mocy twierdzenia
\ref{ma2:thm:series-grouping}. Szereg zgrupowany jest
iloczynem Cauchy'ego względem $T$, a jego suma jest
równa $T(x,y)$ z wniosku \ref{ma2:cor:bilinear-absolute}.
Każde inne wyliczenie par jest przestawieniem
pierwszego; twierdzenie \ref{ma2:thm:absolute-rearrangement}
zapewnia zachowanie sumy.
\end{proof}

W następnym rozdziale przejdziemy do iloczynów nieskończonych,
czyli granic iloczynów częściowych, a później do rodzin
sumowalnych i szeregów wielokrotnych.
